\section{Discussion and Limitations}\label{sec:discussion}

\paragraph{What reconstruction quality does and does not tell us.} Reconstruction quality is informative about
behaviour. Within a GPT-2 layer, SAEs with lower reconstruction error damage next-token prediction less; across depth,
reconstruction and splice damage degrade together; and the large reconstruction gains of the scaled SAEs brought large,
confirmed behavioural gains. It is not informative in the same way about information or causal preservation. It does
not robustly predict truth-information loss within a layer, SAEs that reconstruct well still lose persona information,
the scaled SAEs brought no confirmed gain in truth information, and at layer 11 the architecture with higher explained
variance retains less of the causal effect (an ordering not re-tested at matched task-level L0). A plausible explanation is that reconstruction metrics weight directions by
their variance, whereas task information and causal effects live in small differences between nearly identical
activations, which an SAE can reconstruct poorly even when it reconstructs each activation well.

\paragraph{Measurement choices change information conclusions.} A frozen probe mixes two questions: whether the
information is still present, and whether the reconstruction stays where that probe was fitted. On truth statements
both contribute: much of the frozen drop is generic reader sensitivity, but about a third to a half of it at layers
6--11 is SAE error aligned with the truth labels. The same reader produced the only apparent truth gain of
fine-tuning. Information claims
should therefore name the reader, and comparisons with PCA should name what is matched: equal reconstruction error asks
whether the SAE's error is unusually costly, whereas an equal per-token budget asks whether sparse, input-dependent
selection beats one fixed subspace of the same width.

\paragraph{Faithfulness is not one scalar.} Each improvement method improved what it optimised and little else.
Truth-focused fine-tuning moved a frozen reader; a truth-subspace penalty improved in-distribution truth information
without transfer to other statement families; scaling improved reconstruction and behaviour; behaviour-targeted
training improved behaviour at a cost elsewhere; and an apparent causal gain from output-KL disappeared when sparsity was
fixed. A claim that one SAE is ``more faithful'' than another is only meaningful when it names the level. The Gemma
result extends the behavioural part of this account to a second model and an independently trained SAE family.

\paragraph{What the held-out persona test confirmed.} The single RepE-LM test pass (Section~\ref{sec:h1test}) confirms
that the scaled SAEs reproduce the honest-minus-dishonest difference more faithfully, in all nine cells and robustly to
statement sampling. It confirms the layer-8 causal gain only for TopK, and that confirmation does not survive adding
statement-sampling uncertainty. Because a neutral one-word swap loses at least as large a share of its signal
(Section~\ref{sec:info}), what the SAEs lose on this task is one-word instruction information in general, not
honesty-specific information.

\paragraph{Limitations.}
\begin{itemize}
  \item \textbf{Scope.} Two models and a finite set of SAEs: three in-house architectures at one operating point, one
        data-scaled recipe and one public family in GPT-2, and one public SAE per layer at three layers of Gemma-2-2B.
        Our JumpReLU uses a surrogate gradient and has many dead latents at layer 0, and a modern JumpReLU could not be
        trained to the target sparsity in the scaling experiment.
  \item \textbf{No controlled cross-model comparison.} Gemma was evaluated at two levels only, and the settings
        differ in too many respects to compare the models.
  \item \textbf{Sparsity matching.} Architectures are matched on OWT, not on task prompts; a TopK experiment trained
        at task-matched sparsity ran at validation level, and only 1 of 4 architecture differences persists after BH
        correction, so those differences are not robust. Post-hoc truncation to a common task-level L0 leaves 54
        of 57 scaled-versus-base comparison classes unchanged but changes 13 of 19 orderings of the base architectures
        on task metrics, so architecture rankings on task prompts are not robust (Appendix~\ref{sec:audit}).
  \item \textbf{Causal evidence.} One persona-instruction effect of 0.02--0.04 nats at three sites of one model. The
        patches retain the error term and do not test mediation through SAE latents. A mediated SAE-in-the-loop
        experiment was run on validation data, but its random-decoder control was also classified as mediated, so
        whether the effect is mediated specifically by SAE latents is inconclusive.
  \item \textbf{Evidence status.} The held-out results are the base reconstruction table, the improvement-phase
        comparisons on OWT and truth data, and one pre-registered RepE-LM test evaluation. In that evaluation the
        per-pair persona direction cosine was confirmed in all nine cells; the causal-preservation and persona re-fit
        hypotheses were not confirmed, only their TopK layer-8 cells passed, and both of those cells are fragile once
        statement sampling is included (Section~\ref{sec:h1test}). The capacity-matched PCA test
        confirmation is a test result (P1 confirmed: no trained SAE worse than PCA in any of the 20 cells). Every other
        persona, causal, capacity-matched PCA, truth-subspace and Gemma result is a validation finding. Three seeds share training data, so seed spread
        understates uncertainty, and fine-tuning criteria are applied per cell without multiplicity correction.
  \item \textbf{Inference over seeds, not statements.} The pre-registered comparison rules treat the evaluation set as
        fixed. With statement and document sampling added (Appendix~\ref{sec:audit}), the TopK cells of HT1 and HT3 and
        two layer-1 $\Delta$CE cells of the OWT test pass are no longer supported.
  \item \textbf{Failed negative control.} In the RepE-LM test, test-only neutral swaps gave the persona probe an AUROC
        just above the pre-registered floor at layers 8 and 11, so the probe is not perfectly specific to the persona
        contrast on new frames. The floor and frames were not changed after this was seen.
  \item \textbf{Construct validity of the persona task.} RepE-LM swaps one instruction word, and a neutral one-word
        swap loses as large a share of its signal, so the persona results measure instruction-token retention, not
        honesty.
  \item \textbf{Pooled truth task.} \texttt{cities} and \texttt{neg\_cities} are pooled, which hides truth that a
        linear probe reads within one polarity from layer 1; low pooled AUROC at early layers is not evidence that truth
        is absent there (Appendix~\ref{sec:audit}).
  \item \textbf{R3.} A pre-registered study of how Gemma Scope width and sparsity affect reconstruction and
        behavioural splicing at the M6 layers has completed for 36 SAEs (widths 16k--65k) at \evlabel{validation}
        level (merged in PR \#30). None of our conclusions depends on it.
\end{itemize}
Natural next steps are to repeat the capacity-matched PCA comparison on test data and on the persona task, to run the mediated intervention, and to extend the information and causal levels
to Gemma-2-2B.
